% 三条路径是待检验的机制；仅上路不要求用户偏好或供给发生变化。
\begin{tikzpicture}[
 x=1mm,y=1mm,font=\figurefont\small,
 block/.style={rectangle,draw=BookBlue,fill=BookBlue!5,
   line width=0.85pt,minimum width=34mm,minimum height=13mm,
   align=center,inner sep=2mm},
 endpoint/.style={block,minimum width=22mm,minimum height=17mm},
 flow/.style={-{Latex[length=2.2mm,width=1.5mm]},draw=BookInk,
   line width=1.1pt,rounded corners=1.2mm},
 wire/.style={draw=BookInk,line width=1.1pt,rounded corners=1.2mm},
 loop/.style={flow,draw=BookMuted,dashed,line width=0.85pt}
]
 \path[use as bounding box] (0,-45) rectangle (164,39);
 \node[endpoint] (action) at (12,0) {当前\\展示行动};
 \node[block] (data) at (55,26) {曝光与选择\\进入日志};
 \node[block] (train) at (106,26) {重新训练\\更新模型};
 \node[block,draw=BookAmber,fill=BookAmber!5] (interest) at (55,0)
   {接触与反馈\\改变兴趣或疲劳};
 \node[block,draw=BookAmber,fill=BookAmber!5] (user) at (106,0)
   {用户状态\\与可用历史};
 \node[block,draw=BookTeal,fill=BookTeal!5] (provider) at (55,-26)
   {机会与收益\\影响提供者};
 \node[block,draw=BookTeal,fill=BookTeal!5] (catalog) at (106,-26)
   {创建、修改、退出\\改变后续目录};
 \node[endpoint,draw=BookTeal,fill=BookTeal!5] (next) at (152,0)
   {下一轮\\可执行决策};

 \draw[wire] (action.east) -- (31,0);
 \draw[flow] (31,0) |- (data.west);
 \draw[flow] (31,0) -- (interest.west);
 \draw[flow] (31,0) |- (provider.west);
 \draw[flow] (data.east) -- (train.west);
 \draw[flow] (interest.east) -- (user.west);
 \draw[flow] (provider.east) -- (catalog.west);
 \draw[wire] (train.east) -- (132,26) -- (132,0);
 \draw[flow] (user.east) -- (next.west);
 \draw[wire] (catalog.east) -- (132,-26) -- (132,0);
 \draw[loop] (next.south) -- (152,-41) -- (12,-41) -- (action.south);

 \node[font=\figurefont\footnotesize,text=BookMuted] at (80,37)
   {数据选择循环：潜在偏好可以固定};
 \node[font=\figurefont\footnotesize,text=BookMuted,fill=white,
       inner sep=1mm] at (80,-41) {执行后继续下一轮交互};
\end{tikzpicture}
